\documentclass[11pt]{article}
\usepackage[margin=1in]{geometry}
\usepackage{amsmath,amssymb,amsthm}

\newtheorem{theorem}{Theorem}
\newtheorem{proposition}[theorem]{Proposition}
\newtheorem{definition}[theorem]{Definition}
\newtheorem{remark}[theorem]{Remark}

\title{Kritchman--Raz via Leivant Machines\\in the Guard Tree System}
\author{Notes for the Chwistek project}
\date{}

\begin{document}
\maketitle

\section{Setup}

The Guard tree system has:
\begin{itemize}
\item Terms: $O$ (leaf), $\mathrm{Pair}(a,b)$ (node), $\mathrm{ap1}(f,t)$, $\mathrm{ap2}(g,t,u)$
\item Functors: $I$, $\mathrm{Fst}$, $\mathrm{Snd}$, $\mathrm{Comp}$, $\mathrm{Comp2}$, $\mathrm{Rec}$, $\mathrm{KT}$; and 2-ary: $\mathrm{Pair}$, $\mathrm{Const}$, $\mathrm{Lift}$, $\mathrm{Post}$, $\mathrm{Fan}$, $\mathrm{IfLf}$, $\mathrm{TreeEq}$
\item Derivation rules: computation axioms + structural rules + Schema~F
\item Coding: $\mathrm{code} : \mathrm{Term} \to \mathrm{Tree}$, $\mathrm{codeF1} : \mathrm{Fun1} \to \mathrm{Tree}$
\item Proof checker: $\mathrm{thFunTFor} : \mathrm{Fun1}$ (object-level), $\mathrm{thFun} : \mathrm{Tree} \to \mathrm{Tree}$ (metalevel)
\end{itemize}

The diagonal G\"odel~II proof (\texttt{GodelII.agda}) uses:
\begin{enumerate}
\item \textbf{thm14E}: each derivation $d$ has a proof encoding whose $\mathrm{thFunTFor}$ image equals $\mathrm{reify}(\mathrm{codeEqn}(\mathrm{eq}))$.
\item \textbf{substTFor}: coded substitution, for internalizing the diagonal.
\item \textbf{treeEqSelf}: $\mathrm{TreeEq}(t,t) = O$ via Schema~F.
\end{enumerate}

The obstruction is \textbf{EvalOK}: the condition that $\mathrm{evalCode}(\mathrm{code}(t)) = \mathrm{evalCode}(\mathrm{code}(u))$ is derivable for the equation $t = u$ under consideration. This fails for invalid equations and cannot be internalized uniformly.

\section{Leivant Machines in the Tree System}

Following Nelson's Elements \S17 and Leivant:

\begin{definition}[Register machine]
A \emph{Leivant machine} in the tree system is a $\mathrm{Fun1}$ of the form
$$\mathrm{prog} = \mathrm{Rec}(\mathrm{init}, \mathrm{Post}(f, \mathrm{Post}(\mathrm{Snd}, \mathrm{sndArg})))$$
where $f : \mathrm{Fun1}$ is \emph{Rec-free}: built from $\mathrm{Pair}$, $\mathrm{Fst}$, $\mathrm{Snd}$, $\mathrm{IfLf}$, $\mathrm{KT}$, $\mathrm{Comp}$, $\mathrm{Comp2}$, $\mathrm{Fan}$, $\mathrm{Post}$, $\mathrm{Lift}$ only.
\end{definition}

Key property (formalized in \texttt{Machine.agda}):
$$\mathrm{ap1}(\mathrm{prog})(\mathrm{natCode}(n+1)) = f(\mathrm{ap1}(\mathrm{prog})(\mathrm{natCode}(n)))$$

The one-step function $f$ uses \emph{no recursion}. A single step is bounded-depth composition of projections and conditionals. This is the Leivant \S3.3 property.

\textbf{Example} (copyStep in \texttt{Machine.agda}):
$$\mathrm{copyStep}(\mathrm{Pair}(r_1, r_2)) = \begin{cases}
\mathrm{Pair}(O, r_2) & \text{if } r_1 = O \\
\mathrm{Pair}(\mathrm{Snd}(r_1), \mathrm{Pair}(\mathrm{Fst}(r_1), r_2)) & \text{if } r_1 = \mathrm{Pair}(a,b)
\end{cases}$$

\section{The Kritchman--Raz Argument}

\subsection{Kolmogorov complexity}

\begin{definition}
$K(x) = \min\{|\mathrm{codeF1}(p)| : \exists c.\; \mathrm{ap1}(p, \mathrm{natCode}(c)) = \mathrm{reify}(x)\}$
\end{definition}

The statement ``$K(\xi) > \ell$'' is $\Pi_1$: for all programs $p$ with $|p| \le \ell$ and all clocks $c$, $\mathrm{ap1}(p, \mathrm{natCode}(c)) \ne \mathrm{reify}(\xi)$.

\subsection{The surprise examination}

\begin{enumerate}
\item \textbf{Pigeonhole}: For each $\ell$, there exists $\xi$ with $K(\xi) > \ell$.
  (More trees than programs of bounded size.)
\item \textbf{Chaitin machine} $C_\ell$: searches for a proof of ``$\exists \xi.\; K(\xi) > \ell$'', extracts~$\xi$.
\item $|C_\ell|$ is bounded by a constant (independent of $\ell$, since $\ell$ is encoded as data).
\item If the system proves consistency $\Rightarrow$ the system can prove pigeonhole $\Rightarrow$ $C_\ell$ terminates $\Rightarrow$ $K(\xi) \le |C_\ell| < \ell$ for large $\ell$ $\Rightarrow$ contradiction.
\end{enumerate}

\section{The New Obstruction}

\subsection{Diagonal proof: EvalOK}

In the diagonal proof, the obstruction is \textbf{EvalOK} (formalized in \texttt{NelsonWCBoundary.agda}):
\begin{itemize}
\item evalCode performs tag dispatch (tagO / tagAp1 / passthrough)
\item This dispatch fails for variable-containing terms
\item substTFor (for ruleInst) doesn't commute with evalCode
\item The WC0 fragment (BLevel = 0) avoids this by excluding ruleInst
\end{itemize}

\subsection{KR proof: Computation Bound}

In the KR proof, a \emph{different} obstruction appears. There is no diagonal, so substTFor/evalCode issues don't arise. Instead:

\begin{definition}[Computation Bound Condition]
$\mathrm{CompBound}(\ell, B)$: for all proof-encoding trees $t$ of size $\le B$, if $\mathrm{thFun}(t) = \mathrm{codeEqn}(\mathrm{eq})$ and $\mathrm{eq}$ has the form ``$K(\xi) > \ell$'', then $C_\ell$ terminates within $B$ steps and outputs~$\xi$.
\end{definition}

The system must prove: ``if consistent, then for each $\ell$, $\mathrm{CompBound}(\ell, B(\ell))$ holds for some $B(\ell)$.''

This requires:
\begin{enumerate}
\item The system can enumerate proofs up to size $B$ (tree enumeration --- definable by Rec).
\item The system can run thFunTFor on each candidate (already built).
\item The system can verify that a candidate proof actually proves ``$K(\xi) > \ell$'' (tree comparison --- IfLf + TreeEq).
\item The system can bound $B(\ell)$ --- this is the \textbf{new side condition}.
\end{enumerate}

\subsection{Comparison}

\begin{center}
\begin{tabular}{l|l|l}
& \textbf{Diagonal} & \textbf{KR} \\
\hline
Self-reference & Fixed point (diagonal lemma) & Chaitin machine (search) \\
Key functor & substTFor (coded substitution) & iterate (bounded search) \\
Obstruction & EvalOK (tag dispatch for variables) & CompBound (proof-size bound) \\
WC fragment & BLevel 0 excludes ruleInst & BLevel 0 excludes induction \\
Nature & Algebraic (code structure vs.\ data) & Computational (search termination) \\
\end{tabular}
\end{center}

\section{Connection to BLevel and Aschieri}

\subsection{BLevel hierarchy}

The WC hierarchy stratifies proofs by backtracking level:
\begin{itemize}
\item BLevel 0: no ruleInst, no variables --- ground computation only
\item BLevel $k$: at most $k$ levels of variable instantiation
\end{itemize}

In the diagonal proof, the G\"odel sentence requires BLevel $\ge 1$ (ruleInst for the diagonal).

In the KR proof, the pigeonhole argument requires BLevel $\ge 1$ for a different reason: the induction (Schema~F / ruleF) introduces variables. The KR parameter $\delta$ (incompressibility threshold $\ell$) determines the proof-size bound $B(\ell)$, which in turn determines how deep the BLevel needs to go.

\textbf{Conjecture}: The KR $\delta$ parameter corresponds to BLevel via:
$$\mathrm{BLevel}(\text{pigeonhole proof for } \ell) \sim \log(\ell)$$
because the pigeonhole argument uses $\log(\ell)$ levels of Schema~F induction.

\subsection{Aschieri's backtracking}

Aschieri's game-semantic analysis counts backtracking moves in proof normalization:
\begin{itemize}
\item In the diagonal case: backtracking depth depends on the proof structure (can be unbounded for Schema~F over complex functors).
\item In the KR case: the Chaitin machine's ``backtracking'' is bounded by the number of programs of size $\le \ell$, which is at most $2^\ell$ --- a \emph{concrete} bound.
\end{itemize}

This suggests that the KR obstruction is more amenable to Aschieri-style complexity analysis: the backtracking is bounded by a computable function of $\ell$, whereas the diagonal backtracking depends on the proof itself.

\section{What This Means for Nelson's Program}

Nelson's program asks: can consistency of a fixed system be proved by feasible means?

The diagonal answer: the obstruction is evalCode (non-structural evaluation). Even for a fixed system, the universal evaluator is not feasibly definable.

The KR answer: the obstruction is CompBound (proof-size bounds). For a fixed system, the Chaitin search terminates in bounded time --- but \emph{bounding the bound} requires knowing the proof-size of the pigeonhole argument, which grows with $\ell$.

\textbf{Key insight}: The KR obstruction is \emph{quantitative} where the diagonal obstruction is \emph{qualitative}. The diagonal proof fails because evalCode can't handle variables (a structural impossibility). The KR proof fails because the bound $B(\ell)$ grows faster than the system can prove (a quantitative gap).

This is exactly the Buss--Tao gap: the system can verify each specific $B(\ell)$ (Nelson's success) but cannot prove that $B$ is total (G\"odel's obstruction). In the tree system, this manifests as: each specific pigeonhole instance is provable (by concrete tree enumeration), but the uniform statement ``for all $\ell$, the pigeonhole instance of size $\ell$ is provable'' requires transfinite induction.

\end{document}
